\chapter{Conclusion and Outlook}
\label{chap:main-conclusion}
\label{chap:conclusion}

The central finding is that regularisation must be judged by its statistical
purpose and its numerical reliability. Slope recovery, prediction, and
inference place different demands on a fitted model. Even procedures that
share an exact-arithmetic Krylov space can behave differently when their
coordinates are calculated in finite precision. The controlled experiments
and tablet application make both distinctions visible.

\section{Answers to the research questions}

For slope recovery, the implemented discrepancy rule gives CG--FPLS the
lowest mean ISE in all three simulation models. Its early stopping produces
relatively concentrated slope estimates and suppresses variation in weak
directions. This benefit does not produce a corresponding advantage for
prediction. CG has higher mean MSPE than Arnoldi in every model, most clearly
when several generating directions carry substantial signal. The
bias--variance decomposition explains why: ordinary slope error and
covariance-weighted prediction error need not reward the same restriction.
Even CG's lower slope variance does not imply lower prediction variance.

The comparison between FPCR and FPLS also depends on the model. FPCR orders
directions by predictor variation, whereas FPLS uses response information
to build its search space. Arnoldi and FPCR give very similar prediction
errors when the signal is concentrated in leading directions. Arnoldi's
advantage is clearer in Model~2, where the slope involves more directions.
These results compare the stated complete procedures, including different
tuning rules. They establish no universal ranking of spectral and iterative
regularisation, and selected component counts are not comparable units of
effective degrees of freedom.

The second question concerns numerical representation. At a common valid
dimension, Raw and Arnoldi minimise the same response least-squares
criterion over the same Krylov space. Their simulation results nearly
coincide on the shared subset where the raw fit is numerically stable,
even though each procedure selects its own dimension. Rare raw failures,
however, produce very large errors that strongly affect full-sample means.
Their impact would be obscured by reporting only medians or removing
numerically inconvenient cases. Reorthogonalisation and QR avoid this
failure mechanism in the studied experiments while preserving the intended
fixed-dimension target.

The tablet analysis supports this distinction. Raw and Arnoldi make very
similar predictions, with benchmark RMSEs of $4.607$ and $4.643$,
respectively. Arnoldi uses five well-conditioned directions, while the
selected raw representation has nominal dimension $55$ and is severely
ill-conditioned. Raw's small conditional RMSE advantage therefore does
not establish numerical reliability or a general prediction advantage.
Changing preprocessing can reverse their ordering. The CG prediction cost
is consistent with the simulation findings, but the empirical comparison
does not isolate the causal effect of adding CG iterations.

For inference, the relevant distinction is between regularising the slope
and approximating a sample moment. The sufficient condition
$\sqrt n\|\widehat r-\widehat K\widehat\beta_m\|=o_P(1)$ makes the
CG-based statistic asymptotically equivalent to a direct-moment statistic.
The separate stopping experiment shows that an estimation-oriented variant
can satisfy its discrepancy threshold while leaving too much residual for
this approximation. Late iteration removes the material algorithmic
difference in the studied designs. It does not by itself solve calibration:
the observed rejection rates also depend on finite-sample behaviour,
discretisation, and how critical values are obtained. Neither $m=70$ nor
the simulation's oracle critical values provide a universal prescription
for observed data.

\section{Scope and further work}

The simulation designs isolate specific mechanisms. They use densely
observed curves, Gaussian scores, homoskedastic independent errors, fixed
grids, and a limited set of slopes and covariance spectra. The early-stopping
results depend on the stated variance pilot and discrepancy constants.
Comparisons with other tuning rules, ridge and spline regularisation,
measurement error, irregular observation, and dependent or heavy-tailed
errors would establish how far the findings extend.

The inference study is a controlled benchmark. Its principal critical
values use known continuous-design quantities, whose covariance weights
differ slightly from those induced by the fixed grid. The power study
instead uses simulated finite-sample null quantiles and considers one
direction of departure. The confidence-set illustration restricts candidate
slopes to five coefficients and summarises many sample-specific envelopes;
it is neither a full-slope coverage experiment for that restricted space
nor a simultaneous confidence band for an observed sample. Feasible
plug-in, multiplier, and bootstrap calibration require separate evaluation.

The spectroscopy study uses one response and one data source. Its benchmark
had been inspected during dataset selection, and its bootstrap intervals
condition on the fitted models. It does not measure uncertainty from
development sampling and tuning, or establish performance across independent
tablet batches, laboratories, or acquisition periods. Instrument~2 measures
the same tablets again. Unknown assay units and the absence of a known
physical slope further limit practical and chemical interpretations.
External validation and resampling that respects available batch information
would be the next empirical steps.

The practical implication is to choose regularisation for the question,
retain a numerically reliable representation, and inspect diagnostics
alongside accuracy. A small prediction error alone cannot establish either
accurate slope recovery or valid inference. Keeping these claims separate
allows the strengths and limits of each procedure to be stated clearly.
